% GENERATED FILE. Edit canonical lessons, then run npm run generate:books.
% canonical-source-hash: fnv1a64:c3bf5e97ad8cd60a
% canonical-lessons: SA-C396-dhvanati, SA-C396-darsanam, SA-C396-pratima, SA-C396-garbhagrham, SA-C396-bhaktah

\chapter{At the Shrine}
\label{ch:sa-at-the-shrine}
\hlchaptermodality{396}

\begin{chapteropening}
\textbf{By the end of this chapter:} \emph{I can say what pilgrims find and hear in a shrine.}

Complete the last lesson of chapter 396: I can say what pilgrims find and hear in a shrine.
\end{chapteropening}

\section[\texorpdfstring{\sk{ध्वनति} (dhvanati)}{ध्वनति (dhvanati)}]{\sk{ध्वनति} (dhvanati) — sounds; rings out}

\label{lesson:SA-C396-dhvanati}



\begin{quote}
Before the new one: say the Sanskrit for takes leave, then the Sanskrit for a wandering ascetic.
\end{quote}

\subsection*{You'll want to know: \sk{ध्वनति}}
\textbf{\sk{ध्वनति}} — \emph{dhvanati} — \textquotedblleft{}sounds; rings out\textquotedblright{}.

\begin{quote}
\sk{मन्दिरे} \sk{शंखः} \sk{ध्वनति।}
\end{quote}

\begin{quote}
\emph{mandire śaṅkhaḥ dhvanati.}
\end{quote}

\begin{quote}
A conch sounds in the temple.
\end{quote}

Past (it …): \sk{शंखः} \sk{ध्वनितवान्}, \emph{śaṅkhaḥ dhvanitavān}.

\begin{cousinweb}
\textbf{\sk{ध्वनति}} (\emph{dhvanati}) comes from the root √dhvan, \textquotedblleft{}to sound\textquotedblright{}, which also gives \emph{dhvani}, \textquotedblleft{}a sound\textquotedblright{}; KRM gives its participle \emph{dhvanita} (MW).
\end{cousinweb}

\begin{grammarlens}[title={What you've built}]
One more word for getting around.
\end{grammarlens}

\subsection*{Your turn}
\begin{itemize}
\raggedright
  \item \emph{Say it:} \emph{dhvanati}
  \item \emph{Say it:} \emph{dhvanati}, once more
  \item \emph{Say it:} \emph{mandire śaṅkhaḥ dhvanati.} — A conch sounds in the temple.
  \item \emph{Recall:} say the Sanskrit for takes leave, then the Sanskrit for a wandering ascetic, then say \emph{dhvanati} again
\end{itemize}

\begin{tcolorbox}[breakable,colback=teal!4,colframe=teal!35!black,title={Before you move on}]
What does \sk{ध्वनति} mean? (\textquotedblleft{}Sounds; rings out\textquotedblright{}.) Say it once more.
\end{tcolorbox}
\section[\texorpdfstring{\sk{दर्शनम्} (darśanam)}{दर्शनम् (darśanam)}]{\sk{दर्शनम्} (darśanam) — a sight of the deity; seeing}

\label{lesson:SA-C396-darsanam}



\begin{quote}
Before the new one: say the Sanskrit for a wandering ascetic, then the Sanskrit for rings out.
\end{quote}

\subsection*{You'll want to know: \sk{दर्शनम्}}
\textbf{\sk{दर्शनम्}} — \emph{darśanam} — \textquotedblleft{}a sight of the deity; seeing\textquotedblright{}.

\begin{quote}
\sk{यात्रिकाः} \sk{प्रातः} \sk{मन्दिरे} \sk{दर्शनं} \sk{लभन्ते।}
\end{quote}

\begin{quote}
\emph{yātrikāḥ prātaḥ mandire darśanaṃ labhante.}
\end{quote}

\begin{quote}
In the morning the pilgrims are granted a sight of the deity in the temple.
\end{quote}

A neuter noun.

\begin{cousinweb}
\textbf{\sk{दर्शनम्}} (\emph{darśanam}) is built on the root √dṛś, \textquotedblleft{}to see\textquotedblright{}, with the ending \emph{-ana}; MW gives \emph{darśana} as seeing, sight and an audience, and the sight of a deity at a shrine is that sense in pilgrims' use (MW).
\end{cousinweb}

\begin{grammarlens}[title={What you've built}]
One more word for getting around.
\end{grammarlens}

\subsection*{Your turn}
\begin{itemize}
\raggedright
  \item \emph{Say it:} \emph{darśanam}
  \item \emph{Say it:} \emph{darśanam}, once more
  \item \emph{Say it:} \emph{yātrikāḥ prātaḥ mandire darśanaṃ labhante.} — In the morning the pilgrims are granted a sight of the deity in the temple.
  \item \emph{Recall:} say the Sanskrit for a wandering ascetic, then the Sanskrit for rings out, then say \emph{darśanam} again
\end{itemize}

\begin{tcolorbox}[breakable,colback=teal!4,colframe=teal!35!black,title={Before you move on}]
What does \sk{दर्शनम्} mean? (\textquotedblleft{}A sight of the deity; seeing\textquotedblright{}.) Say it once more.
\end{tcolorbox}
\section[\texorpdfstring{\sk{प्रतिमा} (pratimā)}{प्रतिमा (pratimā)}]{\sk{प्रतिमा} (pratimā) — an image of a deity}

\label{lesson:SA-C396-pratima}



\begin{quote}
Before the new one: say the Sanskrit for rings out, then the Sanskrit for a sight of the deity.
\end{quote}

\subsection*{You'll want to know: \sk{प्रतिमा}}
\textbf{\sk{प्रतिमा}} — \emph{pratimā} — \textquotedblleft{}an image of a deity\textquotedblright{}.

\begin{quote}
\sk{मन्दिरे} \sk{सुन्दरी} \sk{प्रतिमा} \sk{अस्ति।}
\end{quote}

\begin{quote}
\emph{mandire sundarī pratimā asti.}
\end{quote}

\begin{quote}
There is a beautiful image in the temple.
\end{quote}

A feminine noun.

\begin{cousinweb}
\textbf{\sk{प्रतिमा}} (\emph{pratimā}) is built from \emph{prati}, \textquotedblleft{}against, in return\textquotedblright{}, and the root √mā, \textquotedblleft{}to measure\textquotedblright{}: a thing measured against another, a likeness; MW gives \emph{pratimā} as an image, a figure of a god (MW).
\end{cousinweb}

\begin{grammarlens}[title={What you've built}]
One more word for getting around.
\end{grammarlens}

\subsection*{Your turn}
\begin{itemize}
\raggedright
  \item \emph{Say it:} \emph{pratimā}
  \item \emph{Say it:} \emph{pratimā}, once more
  \item \emph{Say it:} \emph{mandire sundarī pratimā asti.} — There is a beautiful image in the temple.
  \item \emph{Recall:} say the Sanskrit for rings out, then the Sanskrit for a sight of the deity, then say \emph{pratimā} again
\end{itemize}

\begin{tcolorbox}[breakable,colback=teal!4,colframe=teal!35!black,title={Before you move on}]
What does \sk{प्रतिमा} mean? (\textquotedblleft{}An image of a deity\textquotedblright{}.) Say it once more.
\end{tcolorbox}
\section[\texorpdfstring{\sk{गर्भगृहम्} (garbhagṛham)}{गर्भगृहम् (garbhagṛham)}]{\sk{गर्भगृहम्} (garbhagṛham) — an inner sanctum; the inner room of a temple}

\label{lesson:SA-C396-garbhagrham}



\begin{quote}
Before the new one: say the Sanskrit for a sight of the deity, then the Sanskrit for an image of a deity.
\end{quote}

\subsection*{You'll want to know: \sk{गर्भगृहम्}}
\textbf{\sk{गर्भगृहम्}} — \emph{garbhagṛham} — \textquotedblleft{}an inner sanctum; the inner room of a temple\textquotedblright{}.

\begin{quote}
\sk{प्रतिमा} \sk{गर्भगृहे} \sk{तिष्ठति।}
\end{quote}

\begin{quote}
\emph{pratimā garbhagṛhe tiṣṭhati.}
\end{quote}

\begin{quote}
The image stands in the inner sanctum.
\end{quote}

A neuter noun.

\begin{cousinweb}
\textbf{\sk{गर्भगृहम्}} (\emph{garbhagṛham}) is built from \emph{garbha}, \textquotedblleft{}the womb, the inside\textquotedblright{}, and \emph{gṛha}, \textquotedblleft{}a house\textquotedblright{}; MW gives the compound as an inner room and as the sanctuary of a temple where the image stands (MW).
\end{cousinweb}

\begin{grammarlens}[title={What you've built}]
One more word for getting around.
\end{grammarlens}

\subsection*{Your turn}
\begin{itemize}
\raggedright
  \item \emph{Say it:} \emph{garbhagṛham}
  \item \emph{Say it:} \emph{garbhagṛham}, once more
  \item \emph{Say it:} \emph{pratimā garbhagṛhe tiṣṭhati.} — The image stands in the inner sanctum.
  \item \emph{Recall:} say the Sanskrit for a sight of the deity, then the Sanskrit for an image of a deity, then say \emph{garbhagṛham} again
\end{itemize}

\begin{tcolorbox}[breakable,colback=teal!4,colframe=teal!35!black,title={Before you move on}]
What does \sk{गर्भगृहम्} mean? (\textquotedblleft{}An inner sanctum; the inner room of a temple\textquotedblright{}.) Say it once more.
\end{tcolorbox}
\section[\texorpdfstring{\sk{भक्तः} (bhaktaḥ)}{भक्तः (bhaktaḥ)}]{\sk{भक्तः} (bhaktaḥ) — a devotee}

\label{lesson:SA-C396-bhaktah}



\begin{quote}
Before the new one: say the Sanskrit for an image of a deity, then the Sanskrit for an inner sanctum.
\end{quote}

\subsection*{You'll want to know: \sk{भक्तः}}
\textbf{\sk{भक्तः}} — \emph{bhaktaḥ} — \textquotedblleft{}a devotee\textquotedblright{}.

\begin{quote}
\sk{भक्तः} \sk{प्रतिमां} \sk{नमति।}
\end{quote}

\begin{quote}
\emph{bhaktaḥ pratimāṃ namati.}
\end{quote}

\begin{quote}
The devotee bows to the image.
\end{quote}

A masculine noun.

\begin{cousinweb}
\textbf{\sk{भक्तः}} (\emph{bhaktaḥ}) comes from the root √bhaj, \textquotedblleft{}to share, to worship\textquotedblright{}, which also gives \emph{bhajati}, \textquotedblleft{}shares, worships\textquotedblright{}: MW gives \emph{bhakta} as devoted or attached to, and as a worshipper (MW).
\end{cousinweb}

\begin{grammarlens}[title={What you've built}]
One more word for getting around.
\end{grammarlens}

\subsection*{Your turn}
\begin{itemize}
\raggedright
  \item \emph{Say it:} \emph{bhaktaḥ}
  \item \emph{Say it:} \emph{bhaktaḥ}, once more
  \item \emph{Say it:} \emph{bhaktaḥ pratimāṃ namati.} — The devotee bows to the image.
  \item \emph{Recall:} say the Sanskrit for an image of a deity, then the Sanskrit for an inner sanctum, then say \emph{bhaktaḥ} again
\end{itemize}

\begin{tcolorbox}[breakable,colback=teal!4,colframe=teal!35!black,title={Before you move on}]
What does \sk{भक्तः} mean? (\textquotedblleft{}A devotee\textquotedblright{}.) Say it once more.
\end{tcolorbox}
